\chapter{Hydrocephalus}
\index{Hydrocephalus}

\section*{Definition and Overview}
Hydrocephalus is a condition in which excess cerebrospinal fluid (CSF) accumulates in the ventricles (fluid-filled spaces) of the brain, increasing pressure inside the skull. It occurs when CSF is overproduced, blocked, or not absorbed properly. Hydrocephalus can be present at birth (congenital) or develop later (acquired), and it can affect children and adults. Without treatment, the raised pressure can damage the brain. Treatment usually involves a shunt, a thin tube that drains excess fluid, or an endoscopic procedure that creates a new drainage route. With prompt treatment, most people do well.

\section*{Epidemiology}
Congenital hydrocephalus affects around 1 in 1,000 births and is one of the most common birth defects of the nervous system. It often accompanies spina bifida and other conditions. Acquired hydrocephalus can occur at any age after head injury, brain infection, bleeding in the brain, or brain tumours. Normal pressure hydrocephalus affects older adults, causing gait problems, memory loss, and urinary incontinence. Hydrocephalus is managed by neurosurgeons, and with treatment, most affected children grow into healthy adults.

\section*{Aetiology and Pathophysiology}
\begin{itemize}
    \item \textbf{CSF circulation:} cerebrospinal fluid is made in the ventricles, circulates around the brain and spinal cord, and is absorbed into the bloodstream
    \item \textbf{Obstruction:} blockage of the narrow passages between ventricles, from congenital narrowing, tumours, or scarring
    \item \textbf{Congenital causes:} spina bifida, aqueduct stenosis, and other malformations
    \item \textbf{Acquired causes:} head injury, bleeding (intraventricular haemorrhage), meningitis, and brain tumours
    \item \textbf{Normal pressure hydrocephalus:} gradual accumulation in older adults, often without an identified cause
    \item \textbf{Pathophysiology:} accumulated fluid enlarges the ventricles and raises pressure, compressing brain tissue
    \item \textbf{In infants:} the skull bones are not yet fused, so the head can enlarge
\end{itemize}

\section*{Clinical Features}
\begin{itemize}
    \item In infants: rapid head growth, a bulging fontanelle, vomiting, poor feeding, and irritability
    \item Sunsetting eyes: the eyes are forced downward
    \item In older children and adults: headache, nausea, vomiting, and blurred vision
    \item Gait problems, balance difficulties, and falls
    \item Cognitive decline, memory problems, and personality change
    \item Urinary incontinence, especially in normal pressure hydrocephalus
    \item Seizures and reduced consciousness in severe cases
\end{itemize}

\section*{Differential Diagnosis}
\begin{itemize}
    \item Other causes of raised intracranial pressure, including brain tumours and infection
    \item Benign enlargement of the subarachnoid spaces in infants, which resolves on its own
    \item Dementia and other causes of cognitive decline in older adults
    \item Gait disorders from other causes
    \item Imaging distinguishes hydrocephalus from these conditions
\end{itemize}

\section*{When to Seek Medical Care}
See a doctor promptly for a baby with a rapidly growing head or vomiting, and for any new headache with vomiting, vision problems, or neurological symptoms. Head injury and infection need assessment.

\textbf{Contact your local emergency services for:}
\begin{itemize}
    \item Sudden severe headache, vomiting, and drowsiness (raised pressure)
    \item Seizures or loss of consciousness
    \item A baby who is very irritable, lethargic, or not feeding
    \item New weakness, vision loss, or difficulty walking
    \item Fever with a headache and stiff neck (possible meningitis)
    \item Symptoms of shunt failure in someone with a shunt: headache, vomiting, or drowsiness
\end{itemize}

\section*{Diagnosis and Investigations}
\begin{itemize}
    \item \textbf{Imaging:} ultrasound in infants, and CT or MRI in older children and adults
    \item \textbf{Head circumference measurement:} serial measurements in infants
    \item \textbf{Neurological examination:} eye movements, reflexes, and gait
    \item \textbf{Lumbar puncture:} measures pressure and samples fluid in some cases
    \item \textbf{Assessment of the cause:} imaging and tests for tumours, bleeding, or infection
    \item \textbf{Monitoring:} repeat imaging to track ventricle size
\end{itemize}

\section*{Management}
\begin{itemize}
    \item \textbf{Ventriculoperitoneal shunt:} a tube from the brain ventricles to the abdomen drains excess fluid; the most common treatment
    \item \textbf{Endoscopic third ventriculostomy:} a keyhole procedure creating a drainage route, for certain types of obstruction
    \item \textbf{Treatment of the cause:} removal of tumours, treatment of infection
    \item \textbf{Shunt care:} lifelong follow-up and monitoring
    \item \textbf{Shunt complications:} blockage and infection are managed urgently
    \item \textbf{Multidisciplinary care:} neurosurgery, neurology, physiotherapy, and educational support
    \item \textbf{Normal pressure hydrocephalus:} shunting improves gait and symptoms in many patients
\end{itemize}

\section*{Complications}
\begin{itemize}
    \item Brain damage and developmental delay if untreated
    \item Shunt blockage, causing symptoms to return rapidly
    \item Shunt infection, with fever and abdominal symptoms
    \item Overdrainage, causing headache or slit ventricles
    \item Seizures and visual problems
    \item Cognitive and motor disability in severe cases
    \item Repeated surgeries over a lifetime
\end{itemize}

\section*{Prognosis and Outcomes}
With early diagnosis and treatment, the outlook for hydrocephalus is good. Most children with shunted hydrocephalus grow up with normal or near-normal intellect, and many lead independent lives. The outcome depends on the cause, the severity of pressure at diagnosis, and the effectiveness of treatment. Shunts require lifelong monitoring, and complications are managed promptly. Normal pressure hydrocephalus responds well to shunting in many older adults.

\section*{Prevention and Self-Care}
\begin{itemize}
    \item Folate supplementation in pregnancy reduces the risk of neural tube defects, including hydrocephalus associated with spina bifida
    \item Attend antenatal scans that can detect hydrocephalus
    \item Protect the head with helmets for sports and seat belts in cars
    \item Treat head injuries and infections promptly
    \item Families with a shunt should know the signs of shunt failure
    \item Attend all neurosurgical follow-up
    \item Access early intervention and educational support for children
    \item Join support organisations for affected families
\end{itemize}

\section*{Sources and Further Reading}
The following reputable sources provide further information for healthcare students and patients seeking reliable, current advice about hydrocephalus.

\begin{itemize}
    \item Hydrocephalus Support Association of Australia. \textit{Information and support.} https://www.hsa.asn.au/
    \item Healthdirect. \textit{Hydrocephalus.} https://www.healthdirect.gov.au/hydrocephalus
    \item Royal Children's Hospital Melbourne. \textit{Hydrocephalus fact sheet.} https://www.rch.org.au/
    \item NHS. \textit{Hydrocephalus: overview and treatment.} https://www.nhs.uk/conditions/hydrocephalus/
    \item Mayo Clinic. \textit{Hydrocephalus: symptoms and causes.} https://www.mayoclinic.org/
    \item Hydrocephalus Association. \textit{Education and support.} https://www.hydroassoc.org/
\end{itemize}
